\subsection{times.var}
\texttt{times.var} configures the date being processed, ephemerides-related timing and the plot time
window. A large part of the file (ephemerides, seasonal plot-limit computation) is scripted and should
not be edited directly.

\textbf{How the date is actually set.} At the top of the file, \texttt{YEAR}/\texttt{MONTH}/\texttt{DAY}
default to the system date (\texttt{\$(date +\%Y)} etc.), but the checked-in file immediately overrides
them with a hand-set date a few lines below (a leftover from the last manual/debug run). \textbf{Because
this codebase has no external process that refreshes \texttt{times.var} before each run}, whichever of
the two is active when \texttt{run.sh} sources this file is what gets used:
\begin{itemize}
  \item to process ``today'' automatically on every scheduled run, comment out (or remove) the
        hand-set override so the system-date lines take effect;
  \item to reprocess a specific past night, or to pin the date for a manual test run, edit the
        hand-set override directly.
\end{itemize}
There is a second file, \texttt{times.var\_sequential}, which is a template of \texttt{times.var} using
the placeholders \texttt{SET\_THE\_DATE\_YEAR}/\texttt{SET\_THE\_DATE\_MONTH}/\texttt{SET\_THE\_DATE\_DAY}/
\texttt{SET\_THE\_UPLOADFIGS} in place of literal values. Nothing in this codebase currently reads or
substitutes into this template -- it is not sourced by \texttt{run.sh}, and no other script in
\texttt{bin/} references it. It is a leftover from a deployment where an external, per-night wrapper
script generated \texttt{times.var} from it; on this single-configuration codebase either ignore it, or
use it as the starting point for your own such wrapper if you want one, substituting the four
placeholders and copying the result over \texttt{times.var} before invoking \texttt{run.sh}.

\subsubsection{Date and time}
\begin{lstlisting}
YEAR=yyyy, MONTH=mm, DAY=dd : date (UT) of the start of the forecast (see above for how
                     this is actually set).
FRCSTDATECOMPUTED=hh: UT hour at which the forecast used for initialization was computed.
                     [current: 12]
FORWARDFRCSTDATEPLUS=hours: how far in the future, from FRCSTDATECOMPUTED, the first
                     forecast file is valid for. [current: 12]
DATEFIGLT=0/1       : if 1, the date used to label figures is the local-time start of the
                     night; if 0, it is the UT date of the end of the night. [current: 1]
HOUR_SHIFT_PLOT=hours: shift of the plot time axis with respect to UT midnight, to align
                     with the simulation start. [current: 0]
UTOFFSET=hours      : local time offset from UT at the site. [current: -7]
HOUR_LIST="string"  : space-separated list of forcing hours; matching initialization
                     files must exist for each. [current: "00 06 12 18"]
TIMESEND=hh         : UT hour by which the init files are expected to have arrived.
                     [current: 18]
TIMELIMIT=seconds   : how long to keep waiting/retrying past TIMESEND before aborting
                     with an error. [current: 10800]
EXITSEC=seconds     : interval between Meso-NH output times; must match OUTTIMESTEP in
                     run.var. [current: 3600]
\end{lstlisting}
(\texttt{EXITMIN} is simply \texttt{EXITSEC/60}, computed automatically.)

\subsubsection{Plot parameters}
\begin{lstlisting}
PLOTSTARTT/PLOTSTARTEXIT, PLOTFIRSTT/PLOTFIRSTEXIT, PLOTSECONDT/PLOTSECONDEXIT,
PLOTENDT/PLOTENDEXIT       : default (non-seasonal) plot time window, in minutes from
                              simulation start and corresponding Meso-NH output index;
                              overwritten by the seasonal computation below when
                              SEASONSPLOT=1 (the default).
FIRSTMONTH=mm, LASTMONTH=mm: calendar months bounding "summer" for the seasonal plot
                              window. [current: 04, 09]
EMISPHERE=1/-1              : 1 for the northern hemisphere, -1 for the southern.
                              [current: 1]
SEASONSPLOT=0/1             : if 1 (default), the plot time window is computed each
                              night from sunset/sunrise/dusk/dawn (via the Sunset.py /
                              Dusk.py / Dawn.py / Sunrise.py helpers) instead of using
                              the fixed values above. [current: 1]
UPLOADFIGS=0/1              : tells backup.sh whether to upload all figures (1) or only
                              the wind-at-ground figures, and only if WINDSPEEDHIGHLIM is
                              exceeded (0) -- see Section~\ref{sec:usage}. A static
                              per-deployment setting here (there is no longer anything
                              that overrides it per run). [current: 1]
\end{lstlisting}
